% 2013 四川卷文科第 19 题：三棱柱 ABC-A1B1C1，D、D1 分别为 BC、B1C1 的中点，P 在 AD 上。
\begin{tikzpicture}[font=\footnotesize, line width=0.45pt, scale=1]
  \coordinate (C)  at (0.80,3.40);
  \coordinate (A)  at (2.60,2.00);
  \coordinate (B)  at (5.60,2.00);
  \coordinate (C1) at (0.80,1.60);
  \coordinate (A1) at (2.60,0.20);
  \coordinate (B1) at (5.60,0.20);
  \coordinate (D)  at (3.20,2.70);
  \coordinate (D1) at (3.20,0.90);
  \coordinate (P)  at (2.87,2.315);
  % 顶面三角形
  \draw (A) -- (B) -- (C) -- cycle;
  % 可见的侧棱与底面棱
  \draw (A) -- (A1);
  \draw (B) -- (B1);
  \draw (C) -- (C1);
  \draw (A1) -- (B1);
  % 平面 A1BC 的辅助线
  \draw (A1) -- (C);
  \draw (A1) -- (B);
  % 被遮挡的棱
  \draw[dashed] (C1) -- (A1);
  \draw[dashed] (C1) -- (B1);
  \draw[dashed] (D)  -- (D1);
  % 点与标注
  \fill (P) circle (1.4pt);
  \node[above left]  at (C)  {$C$};
  \node[below left]  at (2.48,1.86) {$A$};
  \node[right]       at (B)  {$B$};
  \node[left=3pt]    at (C1) {$C_1$};
  \node[below left]  at (A1) {$A_1$};
  \node[below right=1pt] at (B1) {$B_1$};
  \node[above left]  at (D)  {$D$};
  \node[right]       at ($(D1)+(0.10,0)$) {$D_1$};
  \node[left]        at (P)  {$P$};
\end{tikzpicture}
